\documentclass[10pt,a4paper]{amsart}
\usepackage[margin=19mm]{geometry}
\usepackage{amsmath,amssymb,mathtools}
\usepackage[scr=rsfs,cal=euler]{mathalfa}
\usepackage{xfrac}
\usepackage[unicode,hidelinks,bookmarks=false]{hyperref}
\newcommand{\ie}{i.e.\ }
\newcommand{\cf}{cf.\ }
\newcommand{\resp}{respectively\ }
\newcommand{\Subsection}{\subsection}
\providecommand{\nobreakdash}{\nobreak\hbox{-}}
\setlength{\emergencystretch}{2em}
\multlinegap0pt
\usepackage{bm}
\usepackage{bbm}
\usepackage{dsfont}
\usepackage{xspace}
\usepackage{cancel}
\usepackage{mathtools}
\usepackage{amsthm}
\makeatletter
\@ifundefined{problem}{\newtheorem{problem}{Problem}}{}
\@ifundefined{lemma}{\newtheorem{lemma}[problem]{Lemma}}{}
\@ifundefined{theorem}{\newtheorem{theorem}[problem]{Theorem}}{}
\makeatother
\DeclareMathOperator{\HH}{H} % cohomology
\DeclareMathOperator{\ZZ}{Z} % cocycles
\DeclareMathOperator{\aut}{Aut} % Automorphism
\DeclareMathOperator{\dcl}{dcl} % dcl
\DeclareMathOperator{\defin}{def} % for definable Galois cohomology
\DeclareMathOperator{\tp}{tp} % types
\title[The short exact sequence in definable Galois cohomology]{The short exact sequence in definable Galois cohomology}
\author{David Meretzky}
\date{Source: arXiv:2408.04147v3 (CC BY 4.0)}
\begin{document}
\maketitle

\noindent\emph{Source and licence.} The short exact sequence in definable Galois cohomology. Version arXiv:2408.04147v3 (CC BY 4.0), distributed under the Creative Commons Attribution 4.0 International licence, as recorded in the arXiv licence metadata for this exact version and shown on the abstract page https://arxiv.org/abs/2408.04147v3. Full rights notice: licenses/arxiv-2408.04147-CC-BY-4.0.txt.\par\smallskip

\noindent\textit{Editorial scope.} Counted unit: the Theorem whose source label is \texttt{None} in the article (source0.tex lines 399--401, source label \texttt{None}), delivered together with the article's own complete original proof of it (source0.tex lines 403--445, one \texttt{proof} environment of 43 lines). No mathematics is written, repaired or re-derived: statement and proof are the article's own text line for line. Required context retained uncounted: ZZ inj (lines 221--223); pPHS to cocycle (lines 272--274); cocycle to pPHS (lines 290--292). Every definition, display and label of those lines is retained unchanged. Omitted from this package and not redistributed: 1--220, 224--271, 275--289, 293--398, 402--402, 446--518. Nothing inside a retained excerpt is omitted or altered: every retained body is delivered exactly as the article's lines are written, so no omission receipt of any kind accompanies this package. Citations: the retained excerpts carry no citation command at all, so no bibliography entry of the article is transcribed and no reference is redistributed. Reference adaptation: none was needed and none was made; every reference inside the delivered unit resolves inside it, so no unresolved reference or citation is produced. Printed slips kept literal: no printed word, symbol, hypothesis or number of the counted statement or of its proof was altered. Layout and class: the article's own class is \texttt{article} with packages including amsmath, amsthm, amsfonts, amssymb, enumerate; the wrapper is the lane's pinned \texttt{amsart} editorial wrapper at 10pt on A4, which restates the theorem-like environments the retained text needs and adds the article's own macro definitions that the retained text needs (\texttt{\textbackslash HH}, \texttt{\textbackslash ZZ}, \texttt{\textbackslash aut}, \texttt{\textbackslash dcl}, \texttt{\textbackslash defin}, \texttt{\textbackslash tp}), with extra wrapper packages (bm, bbm, dsfont, xspace, cancel, mathtools). The delivered numbering is the unit's own. Assets: no retained span contains a figure environment, a graphics-inclusion command, a PDF-page inclusion, a table or a TikZ picture; no image file is copied into this package, the package declares no asset dependency and no image file is redistributed with it. Licence: version arXiv:2408.04147v3 of this article is distributed under the Creative Commons Attribution 4.0 International licence, as recorded in the arXiv licence metadata for this exact version and shown on the abstract page https://arxiv.org/abs/2408.04147v3. A full rights notice is delivered at licenses/arxiv-2408.04147-CC-BY-4.0.txt. Characters: every retained excerpt is pure ASCII, so the delivered engine, XeLaTeX with its default Latin Modern text font, reports no missing character.
\begin{lemma}\label{ZZ inj}
    Let $\varphi \in \ZZ^1_{\defin}(B/A,G(B))$, then $$- \circ\pi: \ZZ^1_{\defin}(B/A,G(B)) \to \ZZ^1_{\defin}(M/A,G(M))$$ defined by $\varphi \mapsto \varphi \circ \pi$ is an injective map of pointed sets whose image is exactly the set of definable cocycles in $\ZZ^1_{\defin}(M/A,G(M))$ which have definability parameter $\bar{a}$ in $B$. Moreover $-\circ \pi$ descends to an injective map on cohomology, $$\pi^1: \HH^1_{\defin}(B/A,G(B)) \to \HH^1_{\defin}(M/A,G(M)).$$
\end{lemma}

\begin{lemma}\label{pPHS to cocycle}
    An $A$-definable pointed PHS $(X,p)$ for $G$ with $p \in X(B)$ gives rise to a cocycle $\varphi:\aut(M/A) \to G(M)$ which has the $B$-point $p \in X(B)$ as its definability parameter.
\end{lemma}

\begin{lemma}\label{cocycle to pPHS}
    The definable cocycle $\varphi$, with definability parameter contained in $B$, gives rise to a pointed PHS for $G$ whose basepoint is a $B$-point.
\end{lemma}

\begin{theorem}
    The short exact sequence $$1 \to \aut(C/B) \to \aut(C/A) \to \aut(B/A) \to 1$$ induces a short exact sequence in definable Galois cohomology $$1 \to \HH^1_{\defin}(B/A,G(B)) \to \HH^1_{\defin}(C/A,G(C)) \to \HH^1_{\defin}(C/B,G(C))^{\aut(B/A)}.$$
\end{theorem}

\begin{proof}
Let $\varphi$ be a definable cocycle from $\aut(B/A)$ to $G(B)$ with definability data $h$ and $a$. Let $\pi:\aut(C/A) \to \aut(B/A)$ be the quotient map. As in the proof of Lemma \ref{ZZ inj}, we define a map from $Z^1_{\defin}(B/A,G(B))$ to $Z^1_{\defin}(C/A,G(C))$ by sending $\varphi$ to $\varphi \circ \pi$. Now $\varphi \circ \pi$ is still a cocycle witnessed by the same definability data. Since $\pi$ is surjective, this map on cocycles is injective. Let $\varphi_1$ and $\varphi_2$ in  $Z^1_{\defin}(B/A,G(B))$ be cohomologous as witnessed by $g \in G(B)$. As $G(B) \subseteq G(C)$, the same element $g$ witnesses that $\varphi_1 \circ \pi$ and $\varphi_2 \circ \pi$ are cohomologous. So the map on cocycles descends to a well defined map on cohomology which we call $\pi^1$.\\

Let $\varphi_1$ and $\varphi_2$ be cocycles from $\aut(B/A)$ to $G(B)$ with definability data $h_1$, $\bar{a}_1$ and $h_2$, $\bar{a}_2$ respectively. Suppose that $\varphi_1 \circ \pi$ and $\varphi_2\circ\pi$ are cohomologous via $g \in G(C)$. So for all $\sigma \in \aut(C/A)$ and moreover for all $\sigma \in \aut(M/A)$, $$g^{-1}h_1(\bar{a}_1,\sigma(\bar{a}_1))\sigma(g) = h_2(\bar{a}_2,\sigma(\bar{a}_2)).$$ As $\bar{a}_1$ and $\bar{a}_2$ lie in $B$, when $\sigma \in \aut(M/B)$, $h_1(\bar{a}_1,\sigma(\bar{a}_1)) = e$ and $h_2(\bar{a}_2,\sigma(\bar{a}_2)) = e$, and then $g^{-1}\sigma(g) = e$. Thus $g$ is fixed under all automorphisms of $M$ fixing $B$. As $M$ is atomic and strongly $\omega$-homogeneous over $B$, we have $g \in dcl(B)$. As $dcl(B) = B$, we have that $g \in G(B)$. Then $g$ witnesses that $\varphi_1$ and $\varphi_2$ are cohomologous. Thus the induced map on cohomology, $\pi^1$, is injective.\\

Let $\iota:\aut(C/B) \to \aut(C/A)$ be the inclusion map. Precomposition by $\iota$ gives a map of cocycles from $Z^1_{\defin}(C/A,G(C)) \to Z^1_{\defin}(C/B,G(C))$ which induces a well defined map on cohomology $\iota^1$: The same definability data witnesses that precomposition gives a map of definable cocycles and the same element of $G(C)$ witnesses the cohomology relation on the restriction of the cocycle. Note that for cocycles in $\ZZ^1_{\defin}(C/A,G(C))$, precomposition by $\iota$ is equivalent to restriction to $\aut(C/B)$.  \\

We now show that the image of $\iota^1$ is contained in the fixed points of the action of $\aut(B/A)$ on $\HH^1_{\defin}(C/B,G(C))$ via the transform. Let $\sigma \in \aut(C/A)$ be a lift of some $\hat{\sigma} \in \aut(B/A)$. Let $\varphi \in \ZZ^1_{\defin}(C/A,G(C))$, we show that the transform $\sigma(\varphi\circ\iota)$ is cohomologous to $\varphi\circ\iota$. Let $\tau \in \aut(C/B)$. Then
\begin{align*}
    \sigma(\varphi\circ\iota)(\tau) &= \sigma(\varphi(\sigma^{-1}\tau\sigma))\\
    &= \sigma(\varphi(\sigma^{-1}))\sigma(\sigma^{-1}(\varphi(\tau)\tau(\varphi(\sigma))))\\
    &= \sigma(\varphi(\sigma^{-1})\varphi(\tau)\tau(\varphi(\sigma)))\\
    &= g^{-1}\varphi(\tau)\tau(g)
\end{align*}
where $g = \varphi(\sigma) \in G(C)$. So the cocycles $\sigma(\varphi\circ\iota)$ and $\varphi\circ\iota$ are cohomologous via $\varphi(\sigma)$. Hence the image of $\iota^1$ is contained in $\HH^1_{\defin}(C/B,G(C))^{\aut(B/A)}$. Note that we really need a class in the image: If $\varphi$ were an arbitrary cocycle in $\ZZ^1_{\defin}(C/B,G(B))$ then while it is true that $\sigma^{-1}\tau\sigma \in \aut(C/B)$, $\sigma$ is not necessarily in $\aut(C/B)$ and so we could not conclude that $\varphi(\sigma^{-1}\tau\sigma) = \varphi(\sigma^{-1})\sigma^{-1}(\varphi(\tau)\tau(\varphi(\sigma)))$ needed for the second equality of the computation.\\

Choose a cohomology class from $\HH^1_{\defin}(C/A,G(C))$ in the image of $\pi^1$. By definition this class has a representative cocycle $\varphi$ such that $\varphi = \psi \circ \pi$ where $\psi \in Z^1_{\defin}(B/A,G(B))$. Then $\varphi\circ \iota = \psi \circ \pi\circ \iota$ which is the constant cocycle. Since $\iota^1$ is a well defined map on cohomology, the image of the original class must be trivial i.e. cohomologous to the constant cocycle. So the image of $\pi^1$ is contained in the kernel of $\iota^1$. \\

We now check that the kernel of $\iota^1$ is contained in the image of $\pi^1$: Choose a cohomology class from $\HH^1_{\defin}(C/A,G(C))$ in the kernel of $\iota^1$. The first claim is that this class must contain a representative mapping to the constant cocycle: If $\varphi'$ is a representative of a class in the kernel of $\iota^1$ then there is a $g \in G(C)$ such that $\forall \sigma \in \aut(C/B)$, $\varphi'(\sigma) = g^{-1}\sigma(g)$. But then $g\varphi'\sigma(g^{-1})$ is another representative of the class of $\varphi'$ which maps to the constant cocycle.\\

Let $\varphi$ be a representative of our cohomology class in the kernel of $\iota^1$ such that $\varphi \circ \iota$ is the constant cocycle. The cocycle $\varphi$ is a map from $\aut(C/A)$ to $G(C)$ which sends the subgroup $\aut(C/B)$ to the identity of $G(C)$.\\

We now show that the definability data for $\varphi$ can be re-chosen so that the parameter comes from $B$: Let $h(\bar{x},\bar{y})$, $\bar{a}$ be the original definability data for $\varphi$. Let $Z$ be the $A$-definable set $\tp(\bar{a}/A)$; recall that $M$ is isolated over $A$. From $h(\bar{x},\bar{y})$ and $\bar{a}$, we can produce, as in Lemma \ref{cocycle to pPHS}, a pointed PHS for $G$ whose underlying definable set is $(Z\times G)/E$ where $E$ is the $A$-definable equivalence relation defined by $(\bar{a},g)E(\bar{a}',g')$ if $h(\bar{a},\bar{a}')g' = g$. The basepoint of this PHS is the $E$-class $(\bar{a},e)/E$ which furnishes new definability data for $\varphi$ as in Lemma \ref{pPHS to cocycle}.

Now, let $\sigma \in \aut(C/B)$. Then $\sigma((\bar{a},e)/E) = (\bar{a},e)/E$ as $h(\bar{a},\sigma(\bar{a}))e = e$. So $(\bar{a},e)/E$ is fixed pointwise by the action of $\aut(C/B)$, so $(\bar{a},e)/E \in B$.\\

We may now assume that $\varphi$ has definability parameter $\bar{a}$ coming from $B$. As in the proof of Lemma \ref{ZZ inj}, the definability condition $\varphi(\sigma) = h(\bar{a},\sigma(\bar{a}))$ shows $\varphi(\sigma) \in \dcl(\bar{a},\sigma(\bar{a}))$ and so by the normality of $B$ and the fact that $\dcl(B) = B$, the image of $\varphi$ is contained in $G(B)$. As the definability parameter is contained in $B$, $\varphi$ descends to a well defined map $\tilde{\varphi}$ from $\aut(B/A)$ to $G(B)$, and moreover $\tilde{\varphi}\circ\pi = \varphi$. We show that $\tilde{\varphi}$ still satisfies the cocycle condition: Let $\sigma,\tau \in \aut(C/A)$ and $\tilde{\sigma},\tilde{\tau} \in \aut(B/A)$ with $\pi(\sigma) = \tilde{\sigma}$ and $\pi(\tau) = \tilde{\tau}$ 
\begin{align*}
    \tilde{\varphi}(\tilde{\sigma}\tilde{\tau}) &= \varphi(\sigma\tau)\\
    &= \varphi(\sigma)\sigma(\varphi(\tau))\\
    &= \tilde{\varphi}(\tilde{\sigma})\tilde{\sigma}(\tilde{\varphi}(\tilde{\tau}))
\end{align*}
where the last equality holds because the image of $\varphi$ is contained in $B$.  The same definability data for $\varphi$ shows that $\tilde{\varphi}$ is a definable cocycle. The class of $\varphi$ is then the image of the class of $\tilde{\varphi}$ under $\pi^1$.





 



\end{proof}

\end{document}
